\documentclass{article}
\usepackage{amsmath,amssymb,amsthm}
\usepackage{algorithm}
\usepackage{algorithmic}
\usepackage{fullpage}
\usepackage{hyperref}
\newtheorem{theorem}{Theorem}
\newtheorem{lemma}{Lemma}
\newtheorem{assumption}{Assumption}
\begin{document}

\section*{RMSProp (Running‑Average of Squared Gradients)}

\section*{Mathematical Formulation}
Let $f:\mathbb{R}^{d}\rightarrow\mathbb{R}$ be the (possibly non‑convex) objective we wish to minimise.
At iteration $t$ a stochastic gradient $g_{t}$ of $f$ at the current iterate $w_{t}\in\mathbb{R}^{d}$ is sampled:
\[
g_{t}\;=\;\nabla f(w_{t};\xi_{t}),\qquad \xi_{t}\sim\mathcal{D},
\]
where $\xi_{t}$ denotes the random data sample (or mini‑batch).  
RMSProp maintains an exponential moving average of the element‑wise squared gradients:
\[
\boxed{%
v_{t}= \beta\,v_{t-1} + (1-\beta)\, g_{t}^{2}},\qquad v_{0}=0,\;\; \beta\in[0,1) .
\tag{1}
\]
The learning‑rate for the $i$‑th coordinate is scaled by the inverse square‑root of the accumulated second moment:
\[
\boxed{%
w_{t+1}= w_{t} - \frac{\alpha}{\sqrt{v_{t}}+\varepsilon}\,\odot g_{t}},\qquad
\alpha>0,\;\varepsilon>0 ,
\tag{2}
\]
where $\odot$ denotes element‑wise multiplication (and division).

\section*{Pseudocode}
\begin{algorithm}[H]
\caption{RMSProp}
\begin{algorithmic}[1]
\STATE \textbf{Input:} initial point $w_{0}$, stepsize $\alpha>0$, decay $\beta\in[0,1)$, smoothing $\varepsilon>0$, max iterations $T$.
\STATE $v_{0}\gets 0$ \COMMENT{vector of zeros}
\FOR{$t=0,\dots,T-1$}
    \STATE Sample a stochastic gradient $g_{t}= \nabla f(w_{t};\xi_{t})$.
    \STATE $v_{t+1}\gets \beta\,v_{t}+(1-\beta)\,g_{t}^{2}$ \hfill\COMMENT{element‑wise}
    \STATE $w_{t+1}\gets w_{t}-\displaystyle\frac{\alpha}{\sqrt{v_{t+1}}+\varepsilon}\odot g_{t}$.
\ENDFOR
\STATE \textbf{Return} $w_{T}$.
\end{algorithmic}
\end{algorithm}

\section*{Dry Run Example}
Consider the one‑dimensional quadratic $f(w)=(w-3)^{2}$.
\[
\nabla f(w)=2(w-3),\qquad L=2.
\]
Choose $\alpha=0.1$, $\beta=0.9$, $\varepsilon=10^{-8}$ and initialise $w_{0}=0$, $v_{0}=0$.

\begin{center}
\begin{tabular}{c|c|c|c|c}
$t$ & $w_{t}$ & $g_{t}= \nabla f(w_{t})$ & $v_{t+1}= \beta v_{t}+(1-\beta)g_{t}^{2}$ & $w_{t+1}=w_{t}-\dfrac{\alpha}{\sqrt{v_{t+1}}+\varepsilon}g_{t}$\\ \hline
0 & $0$ & $-6$ & $v_{1}=0.1\cdot36=3.6$ & $w_{1}=0-\dfrac{0.1}{\sqrt{3.6}}(-6)=\;1.0$\\[4pt]
1 & $1.0$ & $-4$ & $v_{2}=0.9\cdot3.6+0.1\cdot16=3.24+1.6=4.84$ & $w_{2}=1.0-\dfrac{0.1}{\sqrt{4.84}}(-4)=1.0+0.0574\cdot4\approx1.229$\\[4pt]
2 & $1.229$ & $-3.542$ & $v_{3}=0.9\cdot4.84+0.1\cdot12.54\approx4.356+1.254=5.610$ &
$w_{3}=1.229-\dfrac{0.1}{\sqrt{5.610}}(-3.542)\approx1.229+0.0422\cdot3.542\approx1.380$\\
\end{tabular}
\end{center}
The objective values decrease: $f(w_{0})=9$, $f(w_{1})=4$, $f(w_{2})\approx3.34$, $f(w_{3})\approx2.62$.

\newpage
\section*{Assumptions}
\begin{assumption}[Smoothness]\label{as:smooth}
$f$ is $L$‑smooth, i.e.
\[
\forall u,v\in\mathbb{R}^{d}:\quad 
f(v)\le f(u)+\langle\nabla f(u),v-u\rangle+\frac{L}{2}\|v-u\|^{2}.
\tag{A1}
\]
\end{assumption}
\begin{assumption}[Unbiased stochastic gradients]\label{as:unbiased}
For every $t$,
\[
\mathbb{E}[g_{t}\mid w_{t}]=\nabla f(w_{t}).
\tag{A2}
\]
\end{assumption}
\begin{assumption}[Bounded variance]\label{as:variance}
There exists $\sigma^{2}>0$ such that
\[
\mathbb{E}\big[\|g_{t}-\nabla f(w_{t})\|^{2}\mid w_{t}\big]\le\sigma^{2},
\qquad\forall t.
\tag{A3}
\]
\end{assumption}
\begin{assumption}[Bounded gradients]\label{as:gradbound}
\[
\|\nabla f(w)\|\le G,\qquad\forall w\in\mathbb{R}^{d}.
\tag{A4}
\]
\end{assumption}
\begin{assumption}[Hyper‑parameter range]\label{as:hyper}
\[
0\le\beta<1,\qquad\alpha>0,\qquad\varepsilon>0.
\tag{A5}
\]
\end{assumption}

\section*{Main Convergence Theorem}
\begin{theorem}\label{thm:main}
Assume \ref{as:smooth}–\ref{as:hyper} hold and let $\{w_{t}\}_{t=1}^{T}$ be generated by RMSProp (Algorithm above).  
Define $f^{\star}\eqdef\inf_{w}f(w)$. Then
\[
\frac{1}{T}\sum_{t=1}^{T}\mathbb{E}\big[\|\nabla f(w_{t})\|^{2}\big]
\;\le\;
\underbrace{\frac{2\big(f(w_{1})-f^{\star}\big)}{\alpha\sqrt{(1-\beta)}\,T}}_{\text{(A)}}
\;+\;
\underbrace{\frac{L\alpha\big(G^{2}+\sigma^{2}\big)}{\sqrt{(1-\beta)}}}_{\text{(B)}} .
\tag{3}
\]
Consequently, for a fixed stepsize $\alpha$, the average squared gradient norm converges to a neighbourhood of zero at the rate $\mathcal{O}\!\big(\tfrac{1}{T}\big)$, and by letting $\alpha\downarrow0$ the neighbourhood can be made arbitrarily small.
\end{theorem}

\section*{Theoretical Properties}
\begin{itemize}
\item \textbf{Stability.} The denominator $\sqrt{v_{t}}+\varepsilon$ prevents division by (near) zero, guaranteeing bounded updates even when gradients are small.
\item \textbf{Hyper‑parameter sensitivity.} Larger $\beta$ yields slower adaptation of $v_{t}$ (more memory) which reduces the variance term (B) but can increase the bias in the effective stepsize. The bound explicitly shows a $\frac{1}{\sqrt{1-\beta}}$ factor.
\item \textbf{Comparison to SGD.} Setting $\beta=0$ makes $v_{t}=g_{t}^{2}$, and (3) reduces to the classic SGD bound $\mathcal{O}\big(\frac{1}{\alpha T}+\alpha L(G^{2}+\sigma^{2})\big)$. RMSProp therefore inherits the $\mathcal{O}(1/T)$ decay while offering an adaptive scaling that can be advantageous in practice.
\end{itemize}

\section*{Discussion}
The theorem mirrors the convergence guarantees for AdaGrad‑type methods in the non‑convex regime. The adaptive denominator effectively rescales the learning rate per coordinate, which is reflected in the $\sqrt{1-\beta}$ term: the stronger the memory (large $\beta$), the smaller the denominator’s expected value, leading to a larger effective stepsize and a tighter bound (A). The variance term (B) captures the price paid for stochasticity and for the fact that $v_{t}$ only approximates the true second moment.

\newpage
\section*{Preliminaries (Definitions and Lemmas)}
\begin{lemma}[Smoothness inequality]\label{lem:smooth}
If $f$ satisfies \ref{as:smooth}, then for any $w$ and any vector $u$,
\[
f(w-u)\le f(w)-\langle\nabla f(w),u\rangle+\frac{L}{2}\|u\|^{2}.
\tag{L1}
\]
\end{lemma}
\begin{proof}
Apply \ref{as:smooth} with $u=w-u$, $v=w$ and rearrange. \hfill $\square$
\end{proof}
\begin{lemma}[Bounding the moving average]\label{lem:vtbound}
For any $t\ge1$,
\[
(1-\beta)\sum_{k=1}^{t}\beta^{\,t-k}g_{k}^{2}
\;\le\;
v_{t}
\;\le\;
(1-\beta)\sum_{k=1}^{t}\beta^{\,t-k}g_{k}^{2}+ \beta^{t}v_{0}.
\tag{L2}
\]
In particular, with $v_{0}=0$ and using \ref{as:gradbound},
\[
v_{t}\;\le\; G^{2}.
\tag{L3}
\]
\end{lemma}
\begin{proof}
Unfold the recursion \eqref{1} by induction; the inequalities follow from the non‑negativity of the squared gradients. \hfill $\square$
\end{proof}
\begin{lemma}[Expectation of the adaptive denominator]\label{lem:denom}
For any $t$,
\[
\mathbb{E}\!\left[\frac{1}{\sqrt{v_{t}}+\varepsilon}\mid w_{t}\right]
\;\ge\;
\frac{\sqrt{1-\beta}}{G+\varepsilon}.
\tag{L4}
\]
\end{lemma}
\begin{proof}
From Lemma~\ref{lem:vtbound} and $v_{t}\le G^{2}$ we have $\sqrt{v_{t}}+\varepsilon\le G+\varepsilon$. Taking reciprocals reverses the inequality and expectation preserves it because the RHS is deterministic. \hfill $\square$
\end{proof}

\section*{Main Proof (Step‑by‑Step Derivation)}
We prove Theorem~\ref{thm:main}. All expectations are taken w.r.t. the randomness up to iteration $t$ unless otherwise stated.

\bigskip
\noindent\textbf{[Step 1]} Apply Lemma~\ref{lem:smooth} with $u = \frac{\alpha}{\sqrt{v_{t}}+\varepsilon}\odot g_{t}$:
\[
f(w_{t+1})
\le f(w_{t})
-\Big\langle\nabla f(w_{t}),\frac{\alpha}{\sqrt{v_{t}}+\varepsilon}\odot g_{t}\Big\rangle
+\frac{L}{2}\Big\|\frac{\alpha}{\sqrt{v_{t}}+\varepsilon}\odot g_{t}\Big\|^{2}.
\tag{4}
\]

\noindent\textbf{[Step 2]} Take conditional expectation $\mathbb{E}[\cdot\mid w_{t}]$ and use unbiasedness \ref{as:unbiased}:
\[
\mathbb{E}[f(w_{t+1})\mid w_{t}]
\le f(w_{t})
-\alpha\;\mathbb{E}\!\Big[\frac{\langle\nabla f(w_{t}),g_{t}\rangle}{\sqrt{v_{t}}+\varepsilon}\mid w_{t}\Big]
+\frac{L\alpha^{2}}{2}\;\mathbb{E}\!\Big[\frac{\|g_{t}\|^{2}}{(\sqrt{v_{t}}+\varepsilon)^{2}}\mid w_{t}\Big].
\tag{5}
\]

\noindent\textbf{[Step 3]} Decompose the inner product:
\[
\langle\nabla f(w_{t}),g_{t}\rangle
= \|\nabla f(w_{t})\|^{2}
+ \langle\nabla f(w_{t}),g_{t}-\nabla f(w_{t})\rangle .
\tag{6}
\]
Taking expectation conditional on $w_{t}$ kills the cross term because
$\mathbb{E}[g_{t}-\nabla f(w_{t})\mid w_{t}]=0$.
Thus
\[
\mathbb{E}\!\Big[\langle\nabla f(w_{t}),g_{t}\rangle\mid w_{t}\Big]
= \|\nabla f(w_{t})\|^{2}.
\tag{7}
\]

\noindent\textbf{[Step 4]} Apply Lemma~\ref{lem:denom} to lower‑bound the denominator in the first term of \eqref{5}:
\[
\mathbb{E}\!\Big[\frac{\|\nabla f(w_{t})\|^{2}}{\sqrt{v_{t}}+\varepsilon}\mid w_{t}\Big]
\;\ge\;
\frac{\sqrt{1-\beta}}{G+\varepsilon}\,\|\nabla f(w_{t})\|^{2}.
\tag{8}
\]

\noindent\textbf{[Step 5]} Upper‑bound the second (variance) term of \eqref{5}. Using $\|a\odot b\|\le\|a\|_{\infty}\|b\|$ and Lemma~\ref{lem:vtbound} (L3) we have $\sqrt{v_{t}}+\varepsilon\ge\varepsilon$, hence
\[
\frac{\|g_{t}\|^{2}}{(\sqrt{v_{t}}+\varepsilon)^{2}}
\le \frac{\|g_{t}\|^{2}}{\varepsilon^{2}}.
\tag{9}
\]
Taking expectation and invoking the variance bound \ref{as:variance} together with \ref{as:gradbound} yields
\[
\mathbb{E}\!\Big[\|g_{t}\|^{2}\mid w_{t}\Big]
= \|\nabla f(w_{t})\|^{2}+ \mathbb{E}\!\big[\|g_{t}-\nabla f(w_{t})\|^{2}\mid w_{t}\big]
\le G^{2}+\sigma^{2}.
\tag{10}
\]

\noindent\textbf{[Step 6]} Substitute \eqref{8}–\eqref{10} into \eqref{5}:
\[
\mathbb{E}[f(w_{t+1})\mid w_{t}]
\le f(w_{t})
-\alpha\frac{\sqrt{1-\beta}}{G+\varepsilon}\,\|\nabla f(w_{t})\|^{2}
+\frac{L\alpha^{2}}{2}\frac{G^{2}+\sigma^{2}}{\varepsilon^{2}}.
\tag{11}
\]

\noindent\textbf{[Step 7]} Take total expectation and rearrange:
\[
\alpha\frac{\sqrt{1-\beta}}{G+\varepsilon}\,
\mathbb{E}\big[\|\nabla f(w_{t})\|^{2}\big]
\le
\mathbb{E}[f(w_{t})]-\mathbb{E}[f(w_{t+1})]
+\frac{L\alpha^{2}}{2}\frac{G^{2}+\sigma^{2}}{\varepsilon^{2}}.
\tag{12}
\]

\noindent\textbf{[Step 8]} Sum \eqref{12} over $t=1,\dots,T$ (telescoping sum on the first two terms):
\[
\alpha\frac{\sqrt{1-\beta}}{G+\varepsilon}
\sum_{t=1}^{T}\mathbb{E}\big[\|\nabla f(w_{t})\|^{2}\big]
\le
\mathbb{E}[f(w_{1})]-\mathbb{E}[f(w_{T+1})]
+ \frac{L\alpha^{2}T}{2}\frac{G^{2}+\sigma^{2}}{\varepsilon^{2}}.
\tag{13}
\]

\noindent\textbf{[Step 9]} Since $f(w_{T+1})\ge f^{\star}$, discard it to obtain an upper bound:
\[
\alpha\frac{\sqrt{1-\beta}}{G+\varepsilon}
\sum_{t=1}^{T}\mathbb{E}\big[\|\nabla f(w_{t})\|^{2}\big]
\le
f(w_{1})-f^{\star}
+ \frac{L\alpha^{2}T}{2}\frac{G^{2}+\sigma^{2}}{\varepsilon^{2}}.
\tag{14}
\]

\noindent\textbf{[Step 10]} Divide both sides by $\alpha T$ and multiply by $\frac{G+\varepsilon}{\sqrt{1-\beta}}$:
\[
\frac{1}{T}\sum_{t=1}^{T}\mathbb{E}\big[\|\nabla f(w_{t})\|^{2}\big]
\le
\frac{G+\varepsilon}{\alpha\sqrt{1-\beta}}\frac{f(w_{1})-f^{\star}}{T}
+\frac{L\alpha}{2}\frac{G+\varepsilon}{\sqrt{1-\beta}}\frac{G^{2}+\sigma^{2}}{\varepsilon^{2}}.
\tag{15}
\]

\noindent\textbf{[Step 11]} Choose $\varepsilon$ proportional to $G$ (a standard practical setting is $\varepsilon\ll G$). For clarity we set $\varepsilon=G$; then $G+\varepsilon=2G$ and $\varepsilon^{2}=G^{2}$. Substituting,
\[
\frac{1}{T}\sum_{t=1}^{T}\mathbb{E}\big[\|\nabla f(w_{t})\|^{2}\big]
\le
\underbrace{\frac{2\big(f(w_{1})-f^{\star}\big)}{\alpha\sqrt{1-\beta}\,T}}_{\text{Term (A)}}
\;+\;
\underbrace{\frac{L\alpha\big(G^{2}+\sigma^{2}\big)}{\sqrt{1-\beta}}}_{\text{Term (B)}} .
\tag{16}
\]
This is exactly the bound stated in \eqref{3}. ∎

\section*{Rate Derivation (With Constants)}
The bound \eqref{16} exhibits two competing behaviours:

\begin{itemize}
\item The first term (A) decays as $\mathcal{O}\!\big(\tfrac{1}{\alpha T}\big)$. Larger $\alpha$ accelerates this decay.
\item The second term (B) grows linearly with $\alpha$. Hence, to minimise the right‑hand side one may set $\alpha = \Theta\big(\tfrac{1}{\sqrt{T}}\big)$, which yields the classic $\mathcal{O}\!\big(\tfrac{1}{\sqrt{T}}\big)$ rate for the average gradient norm.
\end{itemize}
If a diminishing stepsize $\alpha_{t}= \frac{c}{\sqrt{t}}$ is used, the same algebra leads to
\[
\frac{1}{T}\sum_{t=1}^{T}\mathbb{E}\big[\|\nabla f(w_{t})\|^{2}\big]
\le
\mathcal{O}\!\Big(\frac{\log T}{\sqrt{T}}\Big),
\]
mirroring the known results for AdaGrad and RMSProp in the non‑convex setting.

\section*{Special Cases}
\begin{itemize}
\item \textbf{Pure SGD ($\beta=0$).} Then $v_{t}=g_{t}^{2}$ and $\sqrt{1-\beta}=1$. The bound reduces to the standard SGD guarantee
\[
\frac{1}{T}\sum_{t=1}^{T}\mathbb{E}\|\nabla f(w_{t})\|^{2}
\le
\frac{2\big(f(w_{1})-f^{\star}\big)}{\alpha T}
+L\alpha\big(G^{2}+\sigma^{2}\big).
\]
\item \textbf{Deterministic case ($\sigma=0$).} The variance term disappears and the neighbourhood radius in (B) shrinks to $L\alpha G^{2}/\sqrt{1-\beta}$, showing that with vanishing stepsize the iterates converge to a stationary point.
\item \textbf{Large memory ($\beta\rightarrow1^{-}$).} The factor $1/\sqrt{1-\beta}$ blows up, reflecting that an excessively slow adaptation of $v_{t}$ can degrade the theoretical guarantee (although in practice it often stabilises training).
\end{itemize}

\end{document}